\section{完整 SIMD 覆盖候选版消融}\label{app:full-simd-ablation}

\subsection{S0、S1、S2 实验设计}

本附录对应第 7.6 节提到的完整 SIMD 候选消融快照（\texttt{Snapshot-E}）。其目的不是重述正文正式方案，
而是说明为何正文保留最小可用版本 \texttt{S1} 作为主结论，同时将完整覆盖候选版 \texttt{S2} 降级为边界比较对象。

三种状态的含义如下：
\begin{itemize}
    \item \texttt{S0}：无 SIMD 基线；
    \item \texttt{S1}：正文正式方案，对应最小可用 LSX 浮点 SIMD 路径；
    \item \texttt{S2}：完整 SIMD 覆盖候选版，补齐更多 lowering 与 idiom。
\end{itemize}

本附录的意义在于把“已成立的最小路径”和“尚待继续扩展的更大覆盖范围”清楚分开，从而避免正文主结论被更复杂、
但证据尚未同样稳固的候选版稀释。

\subsection{全部工作负载与关键工作负载结果}

\begin{figure}[htbp]
    \centering
    \includegraphics[width=0.88\textwidth]{mybenchmark/loongarch64-lsx-full-simd-ablation-geomean.pdf}
    \caption{\texttt{S0/S1/S2} 在不同 workload 集合上的几何平均比较}
    \label{fig:app-lsx-full-simd-ablation-geomean}
\end{figure}

\begin{figure}[htbp]
    \centering
    \includegraphics[width=0.88\textwidth]{mybenchmark/loongarch64-lsx-full-simd-ablation-key-speedup.pdf}
    \caption{\texttt{S0/S1/S2} 在关键 workload 上的相对收益比较}
    \label{fig:app-lsx-full-simd-ablation-key}
\end{figure}

图~\ref{fig:app-lsx-full-simd-ablation-geomean} 表明，在全部 JIT workload 几何平均意义下，\texttt{S1}
仍然优于 \texttt{S2}：\texttt{S1/S0 = 1.0318x}，\texttt{S2/S0 = 1.0252x}，
\texttt{S2/S1 = 0.9936x}。然而，图~\ref{fig:app-lsx-full-simd-ablation-key} 又显示，
\texttt{S2} 在 \texttt{lu.small}、\texttt{h2}、\texttt{avrora} 等关键 workload 上继续超过
\texttt{S1}，说明更完整的后端覆盖并非无效，而是收益更集中。

因此，完整覆盖候选版的正确学术定位应当是“功能覆盖更完整、在关键 workload 上继续获益的候选增强版”，而不是
“已经可以直接替代正文正式方案的最终版本”。这一结论与第 7 章的处理方式完全一致。

\subsection{正式方案与候选增强版的边界}

从论文组织角度看，\texttt{S1} 与 \texttt{S2} 的边界可概括为以下三点。

\begin{enumerate}
    \item \textbf{证据强度不同。} \texttt{S1} 已通过第 7 章的正式对比快照与统计验证快照双重约束；
    \texttt{S2} 目前只适合承担附录中的边界比较角色。
    \item \textbf{研究问题不同。} \texttt{S1} 回答的是“最小可用 LSX 浮点路径是否成立”，而 \texttt{S2}
    回答的是“在更大覆盖范围下是否还能继续收益”。
    \item \textbf{正文责任不同。} \texttt{S1} 承担全文主结论，\texttt{S2} 承担后续扩展方向与边界提醒。
\end{enumerate}

也正因如此，正文选择把 \texttt{S1} 写成主结论，而将 \texttt{S2} 明确收束为附录中的候选增强版。
%---------------------------------------------------------------------------%
